Este informe no se escribió a mano: se generó con Xtal a partir de la carpeta
del proyecto. Las cuatro fuentes de datos entran cada una por su comando y todas
terminan siendo lo mismo --- una \emph{medición}: pares $(x,y)$ con su unidad y
su procedencia --- que después se consolidan en gráficos.

\begin{lstlisting}[language=consola, caption={Las cuatro maneras de conseguir una curva.}, captionpos=b]
# 1. Teorica: se evalua la formula, no hay archivo de entrada.
xtal meas formula --id teorica_mag \
     --expr "-10*math::log10((1-(f/f0)^2)^2 + (f/(q*f0))^2)" \
     --from 20 --to 100000 --points 400 --scale log \
     --const f0=1073.02 --const q=2.0430 --x-unit Hz --y-unit dB

# 2. Simulada: ngspice corre el netlist. Deja magnitud y fase.
xtal circuit import fuentes/filtro-rlc.cir --as filtro-rlc
xtal sim ac filtro-rlc --as sim --node "v(out)" \
     --from 20 --to 100000 --sweep dec --points 30

# 3. Medida: el CSV del instrumento, salteando su metadata.
xtal meas import fuentes/medicion_bode.csv --id medida_mag \
     --skip-rows 6 --x-col 0 --y-col 1 --kind measured --x-unit Hz --y-unit dB

# 4. Importada: un rawfile que ya existia (LTspice o ngspice).
xtal raw import fuentes/variante-q-alto.raw --as q_alto --node "v(out)"
\end{lstlisting}

Un gráfico es una \emph{receta} sobre esas mediciones: las referencia por su
identificador y les aplica una vista. No guarda datos adentro, así que la misma
medición puede aparecer en varios gráficos y cada gráfico puede juntar todas las
que haga falta.

\begin{lstlisting}[language=consola, caption={El Bode de la Figura~\ref{fig:bode}, completo.}, captionpos=b]
xtal plot new bode --kind bode \
     --title "Respuesta en frecuencia del filtro RLC" \
     --x-label "Frecuencia" --y-label "Ganancia"

xtal plot add-series bode --measurement teorica_mag  --panel magnitude
xtal plot add-series bode --measurement sim          --panel magnitude
xtal plot add-series bode --measurement medida_mag   --panel magnitude
xtal plot add-series bode --measurement teorica_fase --panel phase
xtal plot add-series bode --measurement sim_fase     --panel phase
xtal plot add-series bode --measurement medida_fase  --panel phase

xtal run --open
\end{lstlisting}

\noindent
En ningún lado se pidió un color, un tipo de línea ni una escala. El eje
logarítmico sale de haber declarado el gráfico como Bode; el trazo de cada curva
sale de cómo se obtuvo el dato; el color, del rol de la señal. Lo que se declara
es lo que se quiere ver, no cómo dibujarlo.

\vspace{0.6em}
\noindent
El proyecto entero se reconstruye desde cero con un solo comando:

\begin{lstlisting}[language=consola]
./reproducir.sh
\end{lstlisting}

\noindent
que borra todo lo derivado (\texttt{mediciones/}, \texttt{graficos/},
\texttt{esquematicos/}, \texttt{salida/}), lo vuelve a generar desde
\texttt{fuentes/} y compila el PDF. Las únicas entradas son los netlists, los
CSV del instrumento, la imagen, el texto de las secciones y el
\texttt{xtal.toml} que las indexa.
